\subsection*{v6 --- Yukawa Stage D and Calderón interface}

Relative to v5, this version adds two new modules and resolves
the standing review-level critique that "Stage C is data, not
derivation":

\begin{itemize}[leftmargin=1.6em, itemsep=2pt]
\item Added module \texttt{YukawaTrilinearForm.lean} with the
    structure
    \texttt{YukawaTrilinearForm K E := \{$\omega : \Lambda^3 E \to_{\ell K} K$\}}
    and two \texttt{Prop} predicates (originally named
    \texttt{Nondegenerate} and \texttt{Primitive}, renamed in
    v7 to \texttt{ContractionInjective} and
    \texttt{ContractionSurjective} respectively).
\item Constructed the contraction map
    \texttt{contraction} $: E \to_{\ell K} ((\Lambda^2 E) \to_{\ell K} K)$ from $\omega$
    explicitly using
    \texttt{exteriorPower.alternatingMapLinearEquiv} together
    with \texttt{AlternatingMap.curryLeft}. The contraction is
    now a Lean \emph{definition}, not a hypothesis.
\item Proved
    \texttt{contractionEquivOfBoth}: under both predicates the
    contraction upgrades to a linear isomorphism. Hence
    \texttt{toPrimitiveIndecomposableYukawaCoupling}: Stage D
    induces Stage C as a theorem.
\item Proved
    \texttt{finrank\_E\_eq\_three}: Stage D directly implies
    $\dim E = 3$ through the full tower.
\item Added module \texttt{CalderonInterface.lean} with the
    structure \texttt{CalderonCertificate} that wraps Paper 1's
    Calder\'on-induced carrier involution as a Lean datum and
    feeds it into the existing \texttt{CarrierInvolution}
    pipeline. The Paper 1 input is now an explicit
    interface, not a silent consumption.
\item Updated \texttt{AxiomCheck.lean} and
    \texttt{AuditCheck.lean} with eight new theorem names.
    All eight depend only on
    $\{\texttt{propext},\ \texttt{Classical.choice},\
    \texttt{Quot.sound}\}$. \texttt{audit.sh} returns
    \texttt{AUDIT:\ PASS}.
\end{itemize}

\smallskip
The Yukawa interface is now a four-stage tower (under
\texttt{Nontrivial}\,$E_-$ at the Stage D entry point)

\[
\underbrace{(\omega, \texttt{CI}, \texttt{CS})}_{\text{Stage D}}
\;\Rightarrow\;
\underbrace{C_\omega : E_- \simeq (\Lambda^2 E_-)^*}_{\text{Stage C}}
\;\Rightarrow\;
\underbrace{\Lambda^3 E_- \simeq K}_{\text{Stage B}}
\;\Rightarrow\;
\underbrace{E_- \simeq K^3}_{\text{Stage A}}
\;\Rightarrow\;
\dim E_- = 3.
\]

Every step is a Lean theorem. The primary algebraic content
sits at Stage D: an alternating trilinear form plus two
explicit structural predicates.

\subsection*{v5 --- Yukawa Stage C}

Relative to v4, this version closes the Yukawa interface
\emph{within} Lean by introducing a third stage:

\begin{itemize}[leftmargin=1.6em, itemsep=2pt]
\item Added module \texttt{YukawaPrimitive.lean} with the
    structure
    \texttt{PrimitiveIndecomposableYukawaCoupling K E := \{contraction : E $\simeq_{\ell K}$ ($\Lambda^2 E$)\textsuperscript{*}, nontrivial : Nontrivial E\}}.
\item Proved \texttt{finrank\_E\_eq\_three} via the dimension
    chain
    $\dim E \overset{\text{iso}}{=} \dim (\Lambda^2 E)^*
    \overset{\text{dual}}{=} \dim \Lambda^2 E
    \overset{\text{ext.\,pow.}}{=} \binom{\dim E}{2}$
    combined with the elementary fact that $n = n(n-1)/2$ and
    $n \ge 1$ force $n = 3$.
\item Proved the converse \texttt{ofFinrankEqThree} and the
    characterisation \texttt{nonempty\_iff\_finrank\_eq\_three}.
\item Bridged Stage C to Stage B via \texttt{toYukawaTopForm}, and
    transitively to Stage A via \texttt{toYukawaRankCertificate}.
\item Updated \texttt{AxiomCheck.lean} and \texttt{AuditCheck.lean}
    with the five new theorem names. All five depend only on
    $\{\texttt{propext},\ \texttt{Classical.choice},\
    \texttt{Quot.sound}\}$.
\end{itemize}

\smallskip
The Yukawa interface is now structured as a three-stage tower

\[
\underbrace{C_\omega : E_- \simeq (\Lambda^2 E_-)^*}_{\text{Stage C}}
\;\Rightarrow\;
\underbrace{\Lambda^3 E_- \simeq K}_{\text{Stage B}}
\;\Rightarrow\;
\underbrace{E_- \simeq K^3}_{\text{Stage A}}
\;\Rightarrow\;
\dim E_- = 3,
\]

with each step a Lean theorem and a characterisation establishing
the equivalence to $\dim E_- = 3$ at every stage. The
\emph{conceptual} input to Stage C is a contraction iso of $E_-$
with $(\Lambda^2 E_-)^*$, which is the mathematical core of
"primitive indecomposable Yukawa type" with no physics words.

\subsection*{v4 --- Yukawa Stage B}

In v4 the Yukawa rank discharge was promoted from a bare Stage-A
rank witness ($E_- \simeq K^3$, a numerical assumption in
linear-equivalence dress) to a Stage-B algebraic hypothesis
($\Lambda^3 E_- \simeq K$, a structural property of $E_-$).

\begin{itemize}[leftmargin=1.6em, itemsep=2pt]
\item Added module \texttt{YukawaTopForm.lean} with the structure
    \texttt{YukawaTopForm K E := \{topFormEquiv : $\Lambda^3$ E
    $\simeq_{\ell K}$ K\}}.
\item Proved \texttt{YukawaTopForm.finrank\_E\_eq\_three} using
    Mathlib's \texttt{exteriorPower.finrank\_eq} and
    \texttt{Nat.choose\_eq\_one\_iff}.
\item Proved the converse \texttt{ofFinrankEqThree} and the
    characterisation \texttt{nonempty\_iff\_finrank\_eq\_three}.
\item Bridged Stage B to Stage A via
    \texttt{toYukawaRankCertificate}.
\item Updated \texttt{AxiomCheck.lean} and \texttt{AuditCheck.lean}
    with the five new theorem names. All five depend only on
    $\{\texttt{propext},\ \texttt{Classical.choice},\
    \texttt{Quot.sound}\}$.
\end{itemize}

After v4, the conceptual gap shrunk to
\[
\text{"primitive indecomposable Yukawa type"}
\;\Longrightarrow\;
\text{Nonempty}\,(\texttt{YukawaTopForm}\,K\,E_-).
\]
v5 closes this gap with the Stage C contraction iso definition.

\subsection*{v3 --- correcting overclaims}

Relative to v2, v3 corrected three classes of overclaim:

\begin{enumerate}[leftmargin=1.6em, itemsep=3pt]
\item \textbf{Scope of the title and the abstract.} "SM
    Hypercharge Spectrum" $\to$ "$3+2$ Hypercharge Carrier". The
    note no longer suggests it has formalised all SM chiral-
    multiplet hypercharges; only the carrier multiplet
    $\mathrm{diag}(-\tfrac13,-\tfrac13,-\tfrac13,\tfrac12,\tfrac12)$.
\item \textbf{Status of the Stage-A certificates.} "Not
    assumptions" $\to$ "structured Stage-A hypotheses, not bare
    numerical assumptions". The Higgs and Yukawa certificates
    are hypotheses in Lean. Calling them not-assumptions
    misrepresents the type-theoretic status. The Yukawa
    structure is renamed from \texttt{PrimitiveYukawaTypeCertificate}
    to \texttt{YukawaRankCertificate}, accurately describing what
    it actually carries.
\item \textbf{Logical direction of the dependency graph.} The v2
    graph had three incorrect arrows. v3 routes
    \texttt{DeterminantCharacter} $\to$
    \texttt{LatticeRigidityGeneral} (not the reverse); removes
    the spurious \texttt{HiggsIndexCertificate} $\to$
    \texttt{YukawaRankCertificate} edge; and removes the
    \texttt{InvolutionProjectors} $\to$
    \texttt{LatticeRigidityGeneral} edge.
\end{enumerate}

Two further substantive additions:

\begin{itemize}[leftmargin=1.6em, itemsep=1pt]
\item \textbf{Bundled main theorems strengthened.} v3 adds
    \texttt{CarrierPremises.primitive\_pair\_eq\_sm},
    \texttt{CarrierPremises.carrier\_polynomial\_Y}, and the
    top-level \texttt{hypercharge\_carrier\_packet} bundling all
    three conclusions $(q_-,q_+)=(-2,3),\ \tr Y=0,\
    6 Y^2 - Y - \mathbf 1 = \mathbf 0$. The previous version only
    bundled $\tr Y = 0$, which was weaker than the abstract
    claimed.
\item \textbf{Universal-$\lambda$ determinant character.} v3 adds
    \texttt{trace\_zero\_of\_det\_one\_forall\_rat} as the proper
    headline, with \texttt{trace\_zero\_of\_det\_one\_at\_two} kept
    as a single-point helper.
\end{itemize}

The v2-level claim now reduces to its accurate form, refined
for v9:
\[
\boxed{
\begin{aligned}
&\text{Lean verifies the algebraic carrier graph from:}\\
&\quad\text{(i) a Calder\'on \emph{idempotent}
     ($\texttt{CalderonProjector}$: $\pi^2 = \pi$ $\Rightarrow$
     $\varepsilon^2 = 1$),}\\
&\quad\text{(ii) a Higgs Stage-B \emph{polynomial structure}
     ($E_+ \simeq \alpha X + \beta Y$ $\Rightarrow$
     $\dim E_+ = 2$),}\\
&\quad\text{(iii) a Yukawa Stage-D \emph{trilinear form}
     ($\omega : \Lambda^3 E_- \to K$ with
     \texttt{ContractionInjective},}\\
&\quad\text{\phantom{(iii) }\texttt{ContractionSurjective}, and
     \texttt{Nontrivial}\,$E_-$ $\Rightarrow$ $\dim E_- = 3$),}\\
&\quad\text{(iv) a primitive oriented
     determinant-preserving carrier condition,}\\
&\text{down to the $3+2$ hypercharge carrier
     $Y = \mathrm{diag}(-\tfrac13,-\tfrac13,-\tfrac13,\tfrac12,\tfrac12)$}\\
&\text{with $(q_-, q_+) = (-2, 3)$, $\tr Y = 0$,
     and $6 Y^2 - Y - \mathbf 1 = \mathbf 0$.}
\end{aligned}
}
\]
All three rank/involution discharges are now \emph{structural}:
the Calder\'on input is an algebraic idempotent (not a wrapped
involution), the Higgs input is a polynomial-space iso (not a
bare rank witness), and the Yukawa input is an alternating
trilinear form plus two operational predicates (not a packaged
iso).

% ---------------------------------------------------------------- %
\section{Next Steps}
\label{sec:next}
% ---------------------------------------------------------------- %

In rough order of ascending intellectual difficulty:

\begin{enumerate}[leftmargin=1.6em, itemsep=2pt]
\item \textbf{CI workflow.}
    Wire \texttt{audit.sh} into a GitHub Actions workflow that
    refuses any merge with \texttt{AUDIT:\ FAIL}.

\item \textbf{Higgs certificate from $\mathbb P^1$ cohomology.}
    Build a Lean construction of \texttt{HiggsIndexCertificate}
    from the cohomology of $\mathcal O(1)$ on $\mathbb P^1$.

\item \textbf{Higgs scheme-cohomology realisation.}\ The
    Stage B \texttt{HiggsTopForm} now identifies $E_+$ with
    the algebraic shadow of $H^0(\mathbb P^1, \mathcal O(1))$;
    a future module should build the scheme-theoretic
    $H^0(\mathbb P^1_K, \mathcal O(1))$ and exhibit an iso to
    \texttt{DegreeOneTwoVarForm K}. This is Mathlib algebraic
    geometry work and reasonably sized.
\item \textbf{Calderón boundary spectral pair.}\ The Stage A+
    \texttt{CalderonProjector} derives $\varepsilon^2 = 1$
    from an algebraic idempotent, and v10's
    \texttt{BoundaryPolarization} derives the idempotent from
    a complement-pair $H = H_+ \oplus H_-$. What remains is the
    analytic Paper-1 theorem producing the polarisation itself
    from boundary elliptic theory (Calder\'on operator on
    Sobolev spaces). This is a major independent
    formalisation project.
\item \textbf{Yukawa: existence of $\omega$ from TFPT
    \emph{canonically}.}\ The Lean chain now provides full
    existence:
    $\dim E_- = 3 \Rightarrow \exists\,\omega$ with
    \texttt{ContractionInjective} and
    \texttt{ContractionSurjective}
    (\texttt{YukawaStageDExistence}, v10). What remains is
    the TFPT-internal question: does the boundary kernel
    produce \emph{the same canonical} $\omega$ (or one
    equivalent up to the action of the carrier symmetries)?
    The abstract bridge
    \texttt{BoundaryYukawaKernelInterface} (v11) gives a
    Lean-typed target for this future TFPT theorem.

\item \textbf{Derived SM representations.}
    Implement
    \texttt{LinearMap.dualHypercharge},
    \texttt{ExteriorPower.hypercharge},
    \texttt{TensorProduct.hypercharge},
    so that the project can claim "SM hypercharge spectrum"
    rather than "carrier hypercharge".

\item \textbf{$\mathbb Z_6$-quotient theorem.}
    The kernel of the carrier action of
    $SU(3)\times SU(2)\times U(1)$ on $\mathbb K^5$ as $\mathbb Z_6$,
    using \texttt{Mathlib.LinearAlgebra.UnitaryGroup}.

\item \textbf{Mathlib upstream.}
    After (2) and (4), \texttt{Polarization},
    \texttt{InvolutionProjectors.toPolarization}, and the
    Mathlib bridge make a clean PR candidate.
\end{enumerate}

% ---------------------------------------------------------------- %
\section{Result}
% ---------------------------------------------------------------- %

\begin{tfptexports}
\textbf{Lean-verified outputs (v11).}
\begin{itemize}[leftmargin=1.5em, itemsep=0pt]
\item Abstract carrier polynomial:
  \texttt{Polarization.sixY\_carrier\_polynomial}.
\item Boundary involution $\Rightarrow$ projectors:
  \texttt{InvolutionProjectors.CarrierInvolution.toPolarization}.
\item Mathlib bridge:
  \texttt{MathlibBridge.Polarization.toCompleteOrthogonalIdempotents}.
\item General integer rigidity:
  \texttt{LatticeRigidityGeneral.primitive\_trace\_free\_pair\_general}.
\item Structural trace identity:
  \texttt{TraceProjection.trace\_linear\_combination\_of\_idempotents}.
\item Determinant character, universal-$\lambda$ headline:
  \texttt{DeterminantCharacter.trace\_zero\_of\_det\_one\_forall\_rat}.
\item Higgs Stage-A discharge:
  \texttt{HiggsIndexShadow.HiggsIndexCertificate.finrank\_Eplus\_eq\_two}.
\item Yukawa Stage-A discharge (rank witness):
  \texttt{YukawaRank.YukawaRankCertificate.finrank\_Eminus\_eq\_three}.
\item Yukawa Stage-B discharge (algebraic top form):
  \texttt{YukawaTopForm.YukawaTopForm.finrank\_E\_eq\_three}.
\item Yukawa Stage-B characterisation:
  \texttt{YukawaTopForm.YukawaTopForm.nonempty\_iff\_finrank\_eq\_three}.
\item Stage-B $\to$ Stage-A bridge:
  \texttt{YukawaTopForm.YukawaTopForm.toYukawaRankCertificate}.
\item Yukawa Stage-C discharge (contraction iso):
  \texttt{YukawaPrimitive.PrimitiveIndecomposableYukawaCoupling.finrank\_E\_eq\_three}.
\item Yukawa Stage-C characterisation:
  \texttt{YukawaPrimitive.PrimitiveIndecomposableYukawaCoupling.nonempty\_iff\_finrank\_eq\_three}.
\item Stage-C $\to$ Stage-B bridge:
  \texttt{YukawaPrimitive.PrimitiveIndecomposableYukawaCoupling.toYukawaTopForm}.
\item Yukawa Stage-D primary datum and constructed contraction:
  \texttt{YukawaTrilinearForm.YukawaTrilinearForm.contraction}.
\item Stage-D operational predicates:
  \texttt{YukawaTrilinearForm.YukawaTrilinearForm.ContractionInjective} and \texttt{.ContractionSurjective}.
\item Stage-D combined predicate
  (\texttt{PrimitiveIndecomposableYukawaCondition}) and
  rank-3 discharge under \texttt{Nontrivial}\,$E$:
  \texttt{YukawaTrilinearForm.YukawaTrilinearForm.finrank\_E\_eq\_three}
  and \texttt{.finrank\_E\_eq\_three\_of\_condition}.
\item Stage-D transparency form (no \texttt{Nontrivial}, $\dim
  \in \{0,3\}$): \texttt{YukawaTrilinearForm.YukawaTrilinearForm.finrank\_E\_eq\_zero\_or\_three}.
\item Stage-D $\to$ Stage-C bridge:
  \texttt{YukawaTrilinearForm.YukawaTrilinearForm.toPrimitiveIndecomposableYukawaCoupling}.
\item Calderón certificate interface:
  \texttt{CalderonInterface.CalderonCertificate.toCarrierInvolution}.
\item Calderón structural derivation (idempotent $\Rightarrow$
  involution): \texttt{CalderonProjector.CalderonProjector.eps\_sq}.
\item Calderón upstream (v10): boundary polarisation $\Rightarrow$
  idempotent:
  \texttt{BoundaryPolarization.BoundaryPolarization.projector\_isIdempotent}
  and \texttt{toCalderonProjector}.
\item Higgs Stage B discharge (polynomial space $\Rightarrow$
  $\dim E_+ = 2$):
  \texttt{HiggsTopForm.HiggsTopForm.finrank\_Eplus\_eq\_two}.
\item Higgs scheme-cohomology shadow (v10):
  \texttt{HiggsSchemeCohomologyShadow.toDegreeOneTwoVarForm}
  (\texttt{LinearEquiv} between the algebraic shadow and
  \texttt{MvPolynomial.homogeneousSubmodule (Fin 2) K 1}).
\item Yukawa Stage D backward (v10):
  \texttt{YukawaStageDExistence.exists\_yukawaTrilinearForm\_of\_finrank\_eq\_three}
  and the characterisation
  \texttt{stage\_D\_iff\_finrank\_eq\_three} (under
  \texttt{Nontrivial}\,$E$).
\item Boundary polarisation converse (v11):
  \texttt{BoundaryPolarization.ofCarrierInvolution}
  --- the Paper-1 interface is a Lean-proved \emph{algebraic}
  equivalence between the polarisation picture and the
  involution picture (two presentations of the same
  finite-dimensional carrier datum; not the analytic
  construction).
\item Seam-winding interface (v11):
  \texttt{SeamWindingInterface.SeamWindingData.toPrimitiveOrientedDeterminantCarrier}
  and the converse, with the SM specialisation
  \texttt{sm\_pair\_eq}.
\item Boundary Yukawa kernel interface (v11):
  \texttt{BoundaryYukawaKernelInterface.BoundaryYukawaKernel.finrank\_E\_eq\_three},
  \texttt{ofFinrankEqThree}, and the characterisation
  \texttt{nonempty\_iff\_finrank\_eq\_three}. \emph{Typed
  packaging of the Stage D condition --- not a new physical
  derivation. The canonicity of $\omega$ from the TFPT
  boundary kernel remains upstream physics work.}
\item Bundled main theorems:
  \texttt{CarrierPremises.primitive\_pair\_eq\_sm},
  \texttt{CarrierPremises.trace\_Y\_eq\_zero},
  \texttt{CarrierPremises.carrier\_polynomial\_Y},
  \texttt{CarrierPremises.hypercharge\_carrier\_packet}.
\item Concrete carrier polynomial:
  \texttt{Hypercharge.Y\_carrier\_polynomial}.
\end{itemize}

\textbf{Audit outputs.}
$1788$ build jobs green; zero \texttt{sorry}; zero
\texttt{admit}; zero \texttt{unsafe} / \texttt{opaque} /
\texttt{partial}; zero domain-specific axioms; every headline
theorem depends only on
$\{\texttt{propext},\ \texttt{Classical.choice},\ \texttt{Quot.sound}\}$.
\end{tfptexports}

% ---------------------------------------------------------------- %
\appendix
\section{Reproducing the Build and Audit}
% ---------------------------------------------------------------- %

\begin{lstlisting}[language=bash]
# 1. Install elan (Lean toolchain manager)
curl -sSf https://raw.githubusercontent.com/leanprover/elan/master/elan-init.sh \
  | sh -s -- -y --default-toolchain none
export PATH="$HOME/.elan/bin:$PATH"

# 2. Enter the project
cd experiments/lean4-carrier-rigidity

# 3. Fetch Mathlib and its pre-built oleans
lake update                  # ~30 s
lake exe cache get           # ~3 min, ~1 GB

# 4. Build the project
lake build                   # ~8 s on top of cache; 1788 jobs

# 5. Run the seven-check audit gate
./scripts/audit.sh
\end{lstlisting}

Expected trailer:

\begin{lstlisting}
=== Final verdict ===
  AUDIT: PASS
\end{lstlisting}

\end{document}
